\documentclass[a4j,12pt]{jarticle}
\pagestyle{empty}

\title{ルジャンドルの多項式の\\一般項からの計算について}
\date{2014年6月8日}
\author{情報工学科 \quad 篠埜 功}

\begin{document}
\maketitle

ルジャンドルの多項式の一般項からの計算について、
4次の場合($P_4(x)$)の計算過程を示す。

ルジャンドルの多項式の一般形は
\[
P_n(x) = \frac{1}{2^n n!}\frac{{\rm d}^n}{{\rm d}x^n}(x^2-1)^n
\]
である。

以下に2通りの計算方法を示す。
1つは展開してから微分を4回行う方法（計算1）、もう1つは微分しながら
展開を行う方法（計算2）である。
計算1の方が見通しが良い。

(計算1)
\begin{eqnarray*}
P_4(x) &=& \frac{1}{2^4 \cdot 4!} \frac{{\rm d}^4}{{\rm d}x^4} (x^2 -1)^4\\
       &=& \frac{1}{2^4 \cdot 4!} \frac{{\rm d}^4}{{\rm d}x^4} (x^4 -2x^2 + 1)^2\\
       &=& \frac{1}{2^4 \cdot 4!} \frac{{\rm d}^4}{{\rm d}x^4} (x^8 - 4x^6 + 6x^4 -4x^2 +1)\\
       &=& \frac{1}{2^4 \cdot 4!} \frac{{\rm d}^3}{{\rm d}x^3} 
                  (8x^7 - 4 \cdot 6x^5 + 6 \cdot 4x^3 -4 \cdot 2 x)\\
       &=& \frac{1}{2^4 \cdot 4!} \frac{{\rm d}^2}{{\rm d}x^2} 
                  (8 \cdot 7x^6 - 4 \cdot 6 \cdot 5 x^4 + 6 \cdot 4 \cdot 3 x^2 -4 \cdot 2)\\
       &=& \frac{1}{2^4 \cdot 4!} \frac{{\rm d}}{{\rm d}x} 
                  (8 \cdot 7 \cdot 6 x^5 - 4 \cdot 6 \cdot 5 \cdot 4 x^3 + 6 \cdot 4 \cdot 3 \cdot 2 x)\\
       &=& \frac{1}{2^4 \cdot 4!} 
                  (8 \cdot 7 \cdot 6 \cdot 5 x^4 - 4 \cdot 6 \cdot 5 \cdot 4 \cdot 3 x^2 + 
                        6 \cdot 4 \cdot 3 \cdot 2)\\
       &=& \frac{1}{8}
                  (35x^4 - 30x^2 + 3)
\end{eqnarray*}

(計算2)
\begin{eqnarray*}
P_4(x) &=& \frac{1}{2^4 \cdot 4!} \frac{{\rm d}^4}{{\rm d}x^4} (x^2 -1)^4\\
       &=& \frac{1}{2^4 \cdot 4!} \frac{{\rm d}^3}{{\rm d}x^3} (4(x^2 -1)^3(2x))\\
       &=& \frac{1}{2 \cdot 4!} \frac{{\rm d}^3}{{\rm d}x^3} (x(x^2 -1)^3)\\
       &=& \frac{1}{2 \cdot 4!} \frac{{\rm d}^2}{{\rm d}x^2} ((x^2 -1)^3 + x \cdot 3(x^2-1)^2(2x))\\
       &=& \frac{1}{2 \cdot 4!} \frac{{\rm d}^2}{{\rm d}x^2} ((x^2 -1)^3 + 6x^2 (x^2-1)^2)\\
       &=& \frac{1}{2 \cdot 4!} \frac{{\rm d}}{{\rm d}x} (3(x^2-1)^2(2x) + 12x(x^2-1)^2 + 6x^2 \cdot 2(x^2-1)(2x))\\
       &=& \frac{6}{2 \cdot 4!} \frac{{\rm d}}{{\rm d}x} 
               (x(x^2-1)^2 + 2x(x^2-1)^2 + 4 x^3 (x^2-1))\\
       &=& \frac{1}{8}
               ((x^2-1)^2 + x \cdot 2(x^2-1)(2x)\\
       &&         + 2(x^2-1)^2 + 2x \cdot 2 (x^2-1)(2x)
                + 12x^2 (x^2-1) + 4x^3 (2x))\\
       &=& \frac{1}{8}
               (35x^4 - 30x^2 + 3)\\
\end{eqnarray*}

なお、この計算は$P_n(x)$の一般項の意味を理解するために示したものであり、
後に講義で紹介する予定のシュミットの直交化という方法で計算ができれば十分
である。$P_n(x)$の一般項についてはこの講義の範囲外とする。

\end{document}

